%% kapa: exam=1 round=33 year=2022 period=1 form=A subject=real_estate subject_ko=부동산학원론 seq=20 number=100
%% source: Q-Net 감정평가사 기출문제 (cst00309) artlSeq=5209997 file=raw/1008/5209997/2022년도 제33회 감정평가사 1차시험 1교시(민법,경제학,부동산학).pdf
%% conversion: manual transcription from page image; review=proofread-sonnet; flags=none
\begin{question}{100}
부동산증권에 관한 설명으로 옳지 않은 것은?
\begin{choices}
  \item 한국주택금융공사는 유동화증권의 발행을 통해 자본시장에서 정책모기지 재원을 조달할 수 있다.
  \item 금융기관은 주택저당증권(MBS)을 통해 유동성 위험을 감소시킬 수 있다.
  \item 저당담보부채권(MBB)의 투자자는 채무불이행위험을 부담한다.
  \item 저당이체증권(MPTS)은 지분형 증권이며 유동화기관의 부채로 표기되지 않는다.
  \item 지불이체채권(MPTB)의 투자자는 조기상환위험을 부담한다.
\end{choices}
\end{question}
\answer{3}
